\clearpage
\section{Методология прогнозного анализа}
\label{sec:modeling-methods}

\subsection{Задача и границы первого эксперимента}

Прогнозный анализ проверяет, улучшают ли характеристики команды предсказание последующего статуса компании сверх сведений о наборе YC и возрасте. Рассматриваются 3\,046 компаний, активных 13 июля 2023 года и имеющих известный статус на 1 августа 2026 года. Положительная метка соответствует поглощению или публичному статусу (272 компании), отрицательная --- активности или неактивности (2\,774 компании). Остальные 64 компании исходной когорты не получают искусственных меток и не участвуют в обучении либо расчёте качества.

Прогнозируется условная вероятность выбранного положительного статуса по исходным характеристикам. Оценка качества непосредственно относится к компаниям с известным исходом. Концентрация неизвестных состояний в группе IK12, установленная в EDA, ограничивает перенос результатов на всю исходную когорту. Отбор по наличию наблюдаемого исхода не устраняется перекрёстной проверкой.

Первый эксперимент ограничен годом и сезоном набора, приближённым возрастом, размером команды и числом основателей. География, тематические теги и тексты оставлены для отдельного расширения. Такое ограничение позволяет проверить полезность сведений о команде при одинаковом составе остальных признаков. Протокол, сетки параметров и случайные начальные значения сохранены до запуска обучения; это фиксация вычислительной процедуры после разведочного анализа, а не предварительная регистрация всего исследования.

\subsection{Сравниваемые модели и обработка признаков}

Отправной точкой служит постоянный прогноз: каждой проверяемой компании назначается доля положительных исходов в соответствующей обучающей части. С ним сравниваются три варианта регуляризованной логистической регрессии: по набору; по набору и возрасту; по набору, возрасту и двум характеристикам команды. Для краткости в таблицах и рисунках логистическая регрессия обозначается ЛР.

Логистическая регрессия задаёт вероятность как
\[
    \widehat p(x)=\frac{1}{1+\exp\{-\eta(x)\}},
    \qquad \eta(x)=\beta_0+\beta^{\mathsf T}z(x),
\]
где $z(x)$ --- преобразованные признаки. Для года набора и возраста используются кубические сплайны с четырьмя равномерно расположенными узлами в диапазоне обучающих значений. Это допускает нелинейную зависимость, не отдавая бустингу преимущество только из-за заведомо линейной базовой модели. За пределами обучающего диапазона сплайн продолжается постоянным значением; отдельная оценка экстраполяции на будущие годы не проводится.

Размер команды заменяется на $\ln(1+x)$ с сохранением нулей. Число основателей включается как числовой признак. Пропуски числовых значений заполняются медианами обучающей части; дополнительно включаются индикаторы их отсутствия. Числовые представления масштабируются по обучающим данным. Сезон кодируется отдельными бинарными столбцами, отсутствие сезона --- отдельной категорией. Используется квадратичная регуляризация с выбором параметра $C\in\{0{,}1;1;10\}$, обратного силе штрафа. Индикаторы пропусков не масштабируются.

Второе семейство --- CatBoost, реализация градиентного бустинга над деревьями с поддержкой категориальных признаков~\cite{catboost-paper}. Сравниваются четыре варианта: набор и возраст; те же признаки с двумя характеристиками команды; исключение размера команды; исключение числа основателей. Сезон передаётся как категориальный признак. Числовые пропуски обрабатываются встроенным механизмом CatBoost, без предварительного заполнения медианами. Преобразование размера команды совпадает с используемым в логистической регрессии.

Для CatBoost перебираются глубины деревьев 3 и 5 и число итераций 200 и 500, всего четыре комбинации. Скорость обучения равна 0,03, коэффициент регуляризации листьев --- 5. Число итераций выбирается во внутренней проверке; внешняя часть не используется для ранней остановки. Это ограниченное сравнение заранее заданных конфигураций, а не поиск наилучшего возможного варианта CatBoost.

Ни одна модель не использует идентификатор компании, поздние характеристики или поздний статус в качестве признака. Взвешивание классов и искусственное выравнивание их численности не применяются: основной объект оценки --- вероятность исхода при наблюдаемой частоте классов. Изменение весов потребовало бы отдельно учитывать его влияние на вероятностные прогнозы.

\subsection{Вложенная перекрёстная проверка}

Настройка параметров отделяется от оценки качества вложенной перекрёстной проверкой~\cite{sklearn-nested}. Во внешнем цикле данные делятся на пять частей. Каждая часть по очереди используется только для проверки, остальные четыре --- для обучения и настройки. Внутри внешней обучающей части выполняется трёхчастная проверка. Для каждой конфигурации суммируются потери на внутренних проверочных наблюдениях и делятся на их общее число. Выбирается конфигурация с наименьшей потерей, затем модель заново обучается на всей внешней обучающей части.

Все оцениваемые по данным преобразования логистической регрессии входят в единую последовательность обучения. Их параметры заново определяются в каждой внутренней и внешней обучающей части, без использования проверочных значений. Таким образом, настройка включает и представление признаков, и модель, что необходимо для предотвращения утечки информации~\cite{kapoor2023}.

Основная схема стратифицирована по целевой переменной и повторяется дважды с начальными значениями генератора 42 и 137. В каждом повторе компания получает ровно один внешний проверочный прогноз от модели, не обучавшейся на ней. Метрики сначала рассчитываются по всем проверочным прогнозам одного повтора, затем усредняются между повторами. Усреднение вероятностей двух повторов перед расчётом основной метрики не производится. Все восемь вариантов используют одинаковые внешние и внутренние разбиения.

Дополнительная схема сохраняет компании одного набора YC в одной части. Группа определяется полным исходным кодом набора, а не только годом. Такое разделение применяется и во внешнем, и во внутреннем циклах с попыткой сохранить соотношение классов. Выполняется один повтор с пятью внешними и тремя внутренними частями. Для каждого разбиения проверяются отсутствие общих групп и наличие обоих классов в обучении и проверке.

Вторая схема характеризует переносимость между наборами компаний. Она не воспроизводит прогнозирование из прошлого в новый календарный период: признаки всех компаний относятся к июлю 2023 года, исходы --- к августу 2026 года. Обе оценки являются внутренней ретроспективной проверкой. Поскольку вся когорта уже рассматривалась в EDA, эти результаты также не представляют независимого подтверждения гипотез, возникших после просмотра данных.

\subsection{Метрики и дополнительная полезность признаков}

Основной критерий настройки и сравнения --- средняя логарифмическая потеря:
\[
    L=-\frac{1}{n}\sum_{i=1}^{n}
    \left[y_i\ln\widehat p_i+(1-y_i)\ln(1-\widehat p_i)\right].
\]
Меньшее значение соответствует лучшему вероятностному прогнозу по этому критерию. Дополнительно рассчитываются средняя точность ранжирования положительного класса (Average Precision, AP), площадь под ROC-кривой (ROC-AUC) и ошибка Брайера $n^{-1}\sum_i(y_i-\widehat p_i)^2$. Первые две метрики оценивают ранжирование, последняя --- квадратичную ошибку вероятностей. AP вычисляется как сумма точности, взвешенной приращениями полноты, а не как трапецеидальная площадь под кривой точности и полноты~\cite{sklearn-metrics}.

Общая доля верных классов не используется для выбора модели: постоянное предсказание отрицательного класса уже даёт 91,07\%. Порог классификации не подбирается, поскольку в задаче не заданы стоимости решений. Калибровка описывается отдельным графиком средних вероятностей и наблюдаемых частот в пяти приблизительно равных по численности группах. Этот график служит диагностикой; дополнительное преобразование вероятностей по нему не обучается.

Для оценки вклада признаков используются парные сравнения моделей на одних и тех же компаниях. Добавление команды сравнивается с моделью по набору и возрасту отдельно для каждого семейства. Для CatBoost дополнительно исключается по одному признаку команды с повторной внутренней настройкой оставшейся модели. Поэтому проверяется дополнительная информация при наличии остальных признаков, а не только величина отдельного коэффициента.

Улучшение определяется как $\Delta L=L_{\text{сравнение}}-L_{\text{новая}}$: положительное значение означает уменьшение ошибки. Для основной схемы сначала усредняется разность индивидуальных потерь компании между двумя повторами. Затем выполняются 2\,000 парных повторных выборок компаний с возвращением и фиксированными прогнозами. Приводятся квантили 2,5\% и 97,5\% распределения средней разности.

Эти границы характеризуют чувствительность сравнения к составу компаний при уже полученных прогнозах. Модели при повторном выборе компаний не переобучаются; зависимость ошибок из-за общих обучающих наблюдений и возможная зависимость внутри наборов не учитываются. Поэтому границы не трактуются как полноценные доверительные интервалы качества всей процедуры обучения и не используются для формального подтверждения гипотез. Отдельно рассматривается изменение результата при разделении наборов YC.

\subsection{Объяснение прогнозов с помощью SHAP}

Для CatBoost с набором, возрастом и двумя характеристиками команды рассчитываются значения SHAP. Они раскладывают исходный выход модели на базовое значение и вклады признаков~\cite{catboost-shap}. В бинарной задаче рассматриваемый исходный выход выражен в единицах логарифма шансов:
\[
    \ln\frac{\widehat p_i}{1-\widehat p_i}
    =\phi_{0,i}+\sum_{j=1}^{5}\phi_{ij}.
\]
Для каждой компании берётся модель её внешней обучающей части в первом стратифицированном повторе. Аддитивность разложения проверяется численно. Базовые значения могут различаться между внешними частями, поскольку соответствуют разным обученным моделям.

Обобщённая важность признака представлена средним $|\phi_{ij}|$ по проверочным компаниям. Она показывает величину его вклада в прогнозы выбранного семейства моделей, но не направление связи и не изменение проверочной ошибки после удаления признака. SHAP не оценивает гиперпараметры: их выбирает внутренняя проверка. Значения SHAP также не интерпретируются как причинные эффекты изменения характеристик компании~\cite{shap-causality}. Они используются только для объяснения уже полученных прогнозов, без последующего отбора признаков в этом эксперименте.
\clearpage
